\documentclass[12pt]{letter}

\usepackage{geometry}
\geometry{
	a4paper,
	total={210mm,297mm},
	left=20mm,
	right=20mm,
	top=25mm,
	bottom=20mm,
}

\usepackage[brazilian]{babel}
\usepackage[utf8]{inputenc}
\usepackage[T1]{fontenc}
\usepackage{fancyhdr}
\usepackage{graphicx}

\renewcommand{\headrulewidth}{0pt}
\fancyhead[L]{
	\includegraphics[width=5cm]{logo_on}
}

\begin{document}

\thispagestyle{fancy}

\begin{flushleft}

\textbf{Modelagem e inversão em coordenadas esféricas na gravimetria}\\
\line(1,0){450}
\\\textbf{Estudante:} Leonardo Uieda\\
\line(1,0){450}
\\\textbf{Orientador:} Valéria Cristina Ferreira Barbosa\\
\line(1,0){450}
\\\textbf{Nível:} Doutorado\\
\line(1,0){450}
\\\textbf{Período Previsto de Bolsa de Estudos:} 2011 - 2015\\
\line(1,0){450}
\\\textbf{Período a que se Refere o Relatório:} 2014\\
\line(1,0){450}

\textbf{Resumo}
\\
\vspace{10 mm}

Existem diversos métodos de inversão linear 3D.

A maioria discretiza a Terra em prismas retangulares retos homogêneos.

O problema inverso consiste em estimar a distribuição de densidade dos prismas
que ajuste a anomalia da gravidade medida.

Levar em conta a curvatura da Terra é desejável em estudos de modelagem em
escala regional (centenas de quilômetros) ou global.

Uma abordagem para isso é discretizar a Terra em prismas esféricos, também
chamados de tesseroides.

Assim, métodos existentes de inversão utilizando prismas podem ser adaptados
para usar tesseroides.

Esta adaptação requer:
um método para calcular o efeito gravitacional de um tesseroide com precisão
e um software que implemente tanto o cálculo direto quanto o método de inversão.


O potencial gravitacional de um tesseroide e suas derivadas não possuem
soluções analíticas.

Foram propostos dois métodos para a integração numérica: a expansão do
integrando em série de Taylor e a Quadratura Gauss-Legendre.

A Quadratura Gauss-Legendre é o método de integração numérica que escolhemos.

A quadratura consiste em aproximar a integral por uma soma ponderada do efeito
gravitacional de N massas pontuais.

Essas massas pontuais são localizadas em pontos no interior do tesseroide que
correspondem a raízes de um polinômio de Legendre de grau N.

O número de massas utilizando (N) é chamado de ``ordem'' da quadratura.

A acurácia da integração numérica depende do número de massas utilizadas e da
distância do tesseroide ao ponto de observação.

Quanto mais próximo estiver o ponto de observação, mais massas são necessárias.

Este comportamento pode ser justificado do ponto de vista da teoria da
amostragem.

É conhecido em métodos potenciais que quanto mais próximo for o ponto de observação, maior é a frequência espacial
(número de onda) do campo potencial.

Consequentemente, na integração numérica são necessários mais pontos para discretizar o integrando sem
que ocorra falseamento.



Um método proposto na literatura para aumentar o número de massas pontuais sem
aumentar a ordem da quadratura é dividir o tesseroide em tesseroides menores.

O campo calculado é a soma dos campos dos tesseroides menores, que é calculado com um número fixo de massas.

Esta divisão se repete para cada tesseroide menor até que ponto de observação
esteja distante o suficiente em relação ao tamanho do tesseroide.

Desta forma, o número de divisões, e consequentemente o número de massas
pontuais, é maior nas partes do tesseroide próximas do ponto de observação.

Podemos dizer que a acurácia da integração numérica depende da razão entre a
distância até o ponto de observação e o tamanho do tesseroide.

Trabalhos anteriores sugerem que esta razão deve ser próxima de 1 (um).

Porém, atualmente não há uma análise que relacione quantitativamente
a maior razão distância/tamanho permitida com o erro cometido na integração.



Neste trabalho, investigamos qual é o erro cometido na integração através da
divisão dos tesseroides.

Utilizamos como referência o modelo de uma casca esférica, para a qual existe
uma solução analítica.

Esta casca é discretizada em tesseroides.

Comparamos o potencial gravitacional e suas primeiras e segundas derivadas
gerados pelo modelo de tesseroides com o valor equivalente para a casca esférica.

Este cálculo foi feito para diversos valores da máxima razão distância/tamanho
permitida.

Fizemos um gráfico do erro cometido na integração pela razão
distância/tamanho.

Desta forma, determinamos qual é o valor da razão distância/tamanho que
proporciona um erro máximo aceitável.

Esta análise pode ser feita para o potencial gravitacional e suas derivadas.





Observamos que o valor ideal para a razão distância/tamanho é diferente para o
potencial e suas derivadas segundas.







\vspace{10 mm}
\textbf{Palavras-Chave:}
Inversão. Gravimetria. Coordenadas esféricas. Tesseroide.

\end{flushleft}
\end{document}
